\documentclass[11pt,a4paper]{article}
\usepackage{fontspec,amsmath,amssymb,amsthm,microtype}
\setmainfont{lmroman10-regular.otf}[BoldFont=lmroman10-bold.otf,ItalicFont=lmroman10-italic.otf,BoldItalicFont=lmroman10-bolditalic.otf]
\usepackage[margin=25mm]{geometry}
\usepackage[unicode,hidelinks]{hyperref}
\setlength{\emergencystretch}{3em}
\setlength{\parskip}{0.4em}
\setlength{\parindent}{0pt}
\allowdisplaybreaks[2]
\raggedbottom
\title{\texorpdfstring{An integrated chain rule from spacetime integrability}{An integrated chain rule from spacetime integrability}}
\author{Omer Deniz Akyildiz\\[0.4em]\normalsize Prepared with AI assistance}
\date{}
\begin{document}\maketitle
Draft mathematical exposition; these notes may contain mistakes.


To differentiate entropy, bounded approximations are first integrated in time. A pointwise time derivative together with joint spacetime \(L^1\) suffices for Fubini and the fundamental theorem of calculus; no common spatial majorant for every time derivative is assumed. The resulting spatial integral is absolutely continuous and has the integrated derivative almost everywhere. A dominated-limit theorem then passes truncated integral identities to the actual entropy identity.



For a spacetime-integrable derivative field \(F(t,x)\), Fubini gives an absolutely continuous time primitive \(A(t)=A(a)+\int_a^t\int F(s,x)d\nu(x)ds\). The fundamental theorem of calculus identifies its derivative almost everywhere with the inner integral. If bounded approximations have convergent endpoint values and a common integrable flux majorant, dominated convergence transfers their integrated identities to the limit. This obtains entropy absolute continuity from an integrated identity rather than assuming pointwise differentiability at every time.


\section{\texorpdfstring{An integrated identity from a spacetime derivative}{An integrated identity from a spacetime derivative}}


Assume \(f(a,\cdot)\in L^1(\nu)\), \(\dot f\in L^1((a,b]\times\nu)\), and \(\partial_t f(t,x)=\dot f(t,x)\) for every \(t\in[a,b]\) and almost every \(x\). Fubini gives \(\dot f(\cdot,x)\in L^1(a,b)\) for almost every \(x\). Apply the scalar fundamental theorem to these sections: \(f(b,x)=f(a,x)+\int_a^b\dot f(t,x)dt\).



The time-integrated derivative is integrable in \(x\) by spacetime integrability. The section identity therefore proves terminal integrability of \(f(b,\cdot)\), rather than requiring it separately. Integrate the identity and exchange time and space to obtain \(\int f(b)d\nu=\int f(a)d\nu+\int_a^b\int\dot f(t,x)d\nu(x)dt\).



For almost every spatial point, integrate its time derivative between the endpoints. Spacetime absolute integrability permits Fubini to exchange the two integrations, yielding the integrated identity for the spatial integral.


\section{\texorpdfstring{Absolute continuity from spacetime integrability}{Absolute continuity from spacetime integrability}}


Apply the integrated fundamental theorem on every subinterval \([a,t]\subseteq[a,b]\). Restricting the spacetime measure preserves integrability, giving both integrability of \(f(t,\cdot)\) and \(F(t)=F(a)+\int_a^t q(u)du\), where \(F(t)=\int f(t,x)d\nu(x)\) and \(q(t)=\int\dot f(t,x)d\nu(x)\).



Fubini gives \(q\in L^1(a,b)\). An indefinite integral of an \(L^1\) function is absolutely continuous; adding the constant \(F(a)\) preserves this property. Equality on the interval transfers absolute continuity to \(F\). Thus regularity of the spatial integral is a conclusion of spacetime integrability and pointwise differentiation.



An integral of an \(L^1\) time function is absolutely continuous. The spacetime norm bound makes the spatially integrated derivative such an \(L^1\) function, supplying this regularity without an assumed entropy derivative.


\section{\texorpdfstring{The almost-\allowbreak{}everywhere derivative of the integrated quantity}{The almost-everywhere derivative of the integrated quantity}}


The subinterval argument gives \(F(t)=F(a)+\int_a^t q(u)du\) for every \(t\in[a,b]\), with \(q\in L^1(a,b)\). Lebesgue differentiation supplies derivative \(q(t)\) for the indefinite integral at almost every interior time. For each such \(t\), an open neighborhood stays inside \((a,b)\), where the displayed identity holds. The local equality transfers the derivative to \(F\), yielding \(F'(t)=\int\dot f(t,x)d\nu(x)\) almost everywhere. No pointwise derivative assertion at every time, or at the endpoints, is inferred.



The fundamental theorem for absolutely continuous primitives gives the derivative almost everywhere. The exceptional time set has measure zero, which is sufficient for all later integrating-factor inequalities.


\section{\texorpdfstring{Passing an integrated identity through dominated limits}{Passing an integrated identity through dominated limits}}


Suppose truncated identities have the form \(B_n-A_n=\int q_n d\nu\). Assume the two endpoint sequences converge to their intended values, \(q_n\to q\) almost everywhere, and \(|q_n|\le g\) with \(g\in L^1(\nu)\). Dominated convergence gives convergence of the right-hand integrals, while the endpoint limits give convergence of the left-hand differences. Their equality for every \(n\) and uniqueness of limits prove \(B-A=\int qd\nu\).



In the entropy application, the endpoints are truncated entropy integrals and \(q_n\) are bounded-test fluxes on a fixed time interval. Spatial energy estimates supply the pointwise flux limit and a time-integrable dominating bound. This lemma passes the integrated identities to the limit without assuming in advance that the full entropy has already been differentiated.



Pass the approximating integrated identities to the limit in both endpoints and flux. Endpoint domination and spacetime flux domination are checked separately because they use different reference measures and majorants.


\section{\texorpdfstring{Regularity supplied by the limiting integral identity}{Regularity supplied by the limiting integral identity}}


Assume only \(q\in L^1(a,b)\) and the exact identity \(H(t)=H(a)+\int_a^t q(u)du\) throughout \([a,b]\). The primitive of \(q\) is absolutely continuous; adding \(H(a)\) and using the identity proves that \(H\) is absolutely continuous.



Lebesgue differentiation gives derivative \(q(t)\) for the primitive at almost every interior time. The identity holds in an interior neighborhood of each such time, so the derivative transfers to \(H\). This gives precisely the final entropy comparison inputs: local absolute continuity and the actual flux as derivative almost everywhere.



The limiting identity expresses the function as a constant plus an integral of an \(L^1\) flux. It therefore supplies both local absolute continuity and its almost-everywhere derivative, closing the regularity needed for entropy comparison.


\end{document}
